\documentclass[10pt,a4paper]{amsart}
\usepackage[margin=19mm]{geometry}
\usepackage{amsmath,amssymb,mathtools}
\usepackage[scr=rsfs,cal=euler]{mathalfa}
\usepackage{xfrac}
\usepackage[unicode,hidelinks,bookmarks=false]{hyperref}
\newcommand{\ie}{i.e.\ }
\newcommand{\cf}{cf.\ }
\newcommand{\resp}{respectively\ }
\newcommand{\Subsection}{\subsection}
\providecommand{\nobreakdash}{\nobreak\hbox{-}}
\setlength{\emergencystretch}{2em}
\multlinegap0pt
\usepackage{bm}
\usepackage{bbm}
\usepackage{dsfont}
\usepackage{xspace}
\usepackage{cancel}
\usepackage{mathtools}
\usepackage{amsthm}
\makeatletter
\@ifundefined{equation}{\newtheorem{equation}{Equation}}{}
\@ifundefined{lemma}{\newtheorem{lemma}[equation]{Lemma}}{}
\@ifundefined{theorem}{\newtheorem{theorem}[equation]{Theorem}}{}
\makeatother
\DeclareMathOperator{\sym}{ \sigma\!\!\!\sigma}
\def\N{\mathbb N}
\def\R{\mathbb R}
\def\S{\mathscr S}
\def\opm{\textup{op}_{\textit{M}}}
\title[A Shubin pseudodifferential calculus on asymptotically conic]{A Shubin pseudodifferential calculus on asymptotically conic manifolds}
\author{Thomas Krainer}
\date{Source: arXiv:2408.08169v2 (CC BY 4.0)}
\begin{document}
\maketitle

\noindent\emph{Source and licence.} A Shubin pseudodifferential calculus on asymptotically conic manifolds. Version arXiv:2408.08169v2 (CC BY 4.0), distributed under the Creative Commons Attribution 4.0 International licence, as recorded in the arXiv licence metadata for this exact version and shown on the abstract page https://arxiv.org/abs/2408.08169v2. Full rights notice: licenses/arxiv-2408.08169-CC-BY-4.0.txt.\par\smallskip

\noindent\textit{Editorial scope.} Counted unit: the Theorem whose source label is \texttt{CompositionCalculus} in the article (source0.tex lines 958--972, source label \texttt{CompositionCalculus}), delivered together with the article's own complete original proof of it (source0.tex lines 973--1120, one \texttt{proof} environment of 148 lines). No mathematics is written, repaired or re-derived: statement and proof are the article's own text line for line. Required context retained uncounted: AnalyticSmoothingApprox (lines 395--397); MelLeibniz (lines 553--558); MelLeibnizRem (lines 560--562); ImprovedMappingProps (lines 855--861); AdjointsCalculus (lines 885--895); AdjointAsymptotics (lines 887--889). Every definition, display and label of those lines is retained unchanged. Omitted from this package and not redistributed: 1--394, 398--552, 559--559, 563--854, 862--884, 890--957, 1121--1601. Nothing inside a retained excerpt is omitted or altered: every retained body is delivered exactly as the article's lines are written, so no omission receipt of any kind accompanies this package. Citations: the retained excerpts carry no citation command at all, so no bibliography entry of the article is transcribed and no reference is redistributed. Reference adaptation: none was needed and none was made; every reference inside the delivered unit resolves inside it, so no unresolved reference or citation is produced. Printed slips kept literal: no printed word, symbol, hypothesis or number of the counted statement or of its proof was altered. Layout and class: the article's own class is \texttt{amsart} with packages including amssymb, amsmath, enumerate, mathrsfs; the wrapper is the lane's pinned \texttt{amsart} editorial wrapper at 10pt on A4, which restates the theorem-like environments the retained text needs and adds the article's own macro definitions that the retained text needs (\texttt{\textbackslash N}, \texttt{\textbackslash R}, \texttt{\textbackslash S}, \texttt{\textbackslash opm}, \texttt{\textbackslash sym}), with extra wrapper packages (bm, bbm, dsfont, xspace, cancel, mathtools). The delivered numbering is the unit's own. Assets: no retained span contains a figure environment, a graphics-inclusion command, a PDF-page inclusion, a table or a TikZ picture; no image file is copied into this package, the package declares no asset dependency and no image file is redistributed with it. Licence: version arXiv:2408.08169v2 of this article is distributed under the Creative Commons Attribution 4.0 International licence, as recorded in the arXiv licence metadata for this exact version and shown on the abstract page https://arxiv.org/abs/2408.08169v2. A full rights notice is delivered at licenses/arxiv-2408.08169-CC-BY-4.0.txt. Characters: every retained excerpt is pure ASCII, so the delivered engine, XeLaTeX with its default Latin Modern text font, reports no missing character.
\begin{lemma}\label{AnalyticSmoothingApprox}
Let $A \in M^{\mu;\vec{\ell}}_O(Z;\overline{\R}_+)$. Then there exist $A_j \in M^{-\infty}_O(Z;\overline{\R}_+)$ such that $A_j \to A$ in $M^{\mu';\vec{\ell}}_O(Z;\overline{\R}_+)$ as $j \to \infty$ for every $\mu' > \mu$. The $A_j \in M^{-\infty}_O(Z;\overline{\R}_+)$ can be explicitly constructed to depend linearly and continuously on $A \in M^{\mu;\vec{\ell}}_O(Z;\overline{\R}_+)$.
\end{lemma}

\begin{equation}\label{MelLeibniz}
\begin{aligned}
p_1{\#}p_2(x,\sigma) &= \frac{1}{2\pi}\iint y^{-i\eta}p_1(x,\sigma+\eta)p_2(xy,\sigma)\,\frac{dy}{y}d\eta \\
&= \sum\limits_{k=0}^{N-1}\frac{1}{k!}[\partial_{\sigma}^kp_1](x,\sigma)[(xD_x)^kp_2](x,\sigma) + r_N(x,\sigma)
\end{aligned}
\end{equation}

\begin{equation}\label{MelLeibnizRem}
r_N(x,\sigma) = \frac{1}{2\pi}\!\int_0^1\!\frac{(1-\theta)^{N-1}}{(N-1)!}\iint y^{-i\eta}[\partial_{\sigma}^Np_1](x,\sigma+\theta\eta)[(xD_{x})^Np_2](xy,\sigma)\,\frac{dy}{y}d\eta d\theta.
\end{equation}

\begin{lemma}\label{ImprovedMappingProps}
Let $p(x,\sigma) = a(x,\sigma,x)$ with $a(x,\sigma,\tau) \in C^{\infty}_B(\R_+,M^{\mu;\vec{\ell}}_O(Z;\overline{\R}_+))$ be arbitrary, and let $\mu \leq 0$. For $j \in \N_0$ with $\mu+\ell_3j \leq 0$ we have
$$
\opm(p),\; [\opm(p)]^* : x^{\alpha}H^{s;(\ell_1,\ell_2)}_b(\R_+\times Z) \to x^{\alpha-j}H^{s-(\mu+\ell_3 j);(\ell_1,\ell_2)}_b(\R_+\times Z)
$$
for all $s,\alpha \in \R$.
\end{lemma}

\begin{theorem}\label{AdjointsCalculus}
Let $A = \opm(p) + G \in \Psi^{\mu;\vec{\ell}}_{\textup{(cl)}}(\R_+\times Z)$. Then the formal adjoint of $A$ with respect to the $L^2_b$-inner product belongs to $\Psi^{\mu;\vec{\ell}}_{\textup{(cl)}}(\R_+\times Z)$. More precisely, suppose that $p(x,\sigma) = a(x,\sigma,x)$ with $a(x,\sigma,\tau) \in C^{\infty}_B(\R_+,M^{\mu;\vec{\ell}}_O(Z;\overline{\R}_+))$. Let then $a^{\star}(x,\sigma,\tau)  \in C^{\infty}_B(\R_+,M^{\mu;\vec{\ell}}_O(Z;\overline{\R}_+))$ be arbitrary with asymptotic expansion
\begin{equation}\label{AdjointAsymptotics}
a^{\star}(x,\sigma,\tau) \sim \sum\limits_{k=0}^{\infty}\frac{1}{k!}[(-xD_x - \tau D_{\tau})^k\partial_{\sigma}^k a](x,\overline{\sigma},\tau)^*
\end{equation}
within this operator-valued symbol class, where the $*$ operator that appears in the asymptotic terms in \eqref{AdjointAsymptotics} is the formal adjoint with respect to the $L^2$-inner product on $Z$. Then there exists $G' \in \Psi^{-\infty}_G(\R_+\times Z)$ such that
$$
A^* = \opm(a^{\star}(x,\sigma,x) ) + G' \in \Psi^{\mu;\vec{\ell}}_{\textup{(cl)}}(\R_+\times Z).
$$
In the classical case we find that $\sym_e(A^*) = \overline{\sym_e(A)}$ (or fiberwise adjoints when the operators are acting in sections of the pull-back to $\R_+\times Z$ of a Hermitian vector bundle on $Z$).
\end{theorem}

\begin{theorem}\label{CompositionCalculus}
Let $A_j = \opm(a_j(x,\sigma,x)) + G_j \in \Psi^{\mu_j;\vec{\ell}}_{\textup{(cl)}}(\R_+\times Z)$, $j = 1,2$, with
$$
a_j(x,\sigma,\tau) \in C^{\infty}_B(\R_+,M^{\mu_j;\vec{\ell}}_{O,\textup{(cl)}}(Z;\overline{\R}_+)).
$$
Then the composition $A_1 A_2 \in \Psi^{\mu_1+\mu_2;\vec{\ell}}_{\textup{(cl)}}(\R_+\times Z)$. More precisely, if $a_1{\#}a_2 \in C^{\infty}_B(\R_+,M^{\mu_1+\mu_2;\vec{\ell}}_{O,\textup{(cl)}}(Z;\overline{\R}_+))$ has the asymptotic expansion
\begin{equation}\label{LeibnizAsymptotics}
(a_1{\#}a_2)(x,\sigma,\tau) \sim \sum\limits_{k=0}^{\infty}\frac{1}{k!}[\partial_{\sigma}^ka_1](x,\sigma,\tau)[(xD_x + \tau D_{\tau})^ka_2](x,\sigma,\tau),
\end{equation}
then there exists $G \in \Psi^{-\infty}_G(\R_+\times Z)$ such that
$$
A_1 A_2 = \opm((a_1{\#}a_2)(x,\sigma,x)) + G.
$$
In the classical case we have $\sym_e(A_1 A_2) = \sym_e(A_1)\sym_e(A_2)$.
\end{theorem}

\begin{proof}
The composition $G_1G_2 \in \Psi^{-\infty}_G(\R_+\times Z)$ in view of the defining mapping properties of the residual class. We next prove that the compositions
$$
\opm(a_1(x,\sigma,x))G_2 \textup{ and } G_1\opm(a_2(x,\sigma,x))
$$
are residual as well. Let $\omega \in C_c^{\infty}(\overline{\R}_+)$ with $\omega \equiv 1$ near $x=0$, and write
$$
a_j(x,\sigma,\tau) = \omega(x) a_j(x,\sigma,\tau) + (1-\omega)(x) a_j(x,\sigma,\tau) = a_{j,0}(x,\sigma,\tau) + a_{j,\infty}(x,\sigma,\tau).
$$
For $K > 0$ large enough we have
$$
a_{j,0}(x,\sigma,x),\, x^{-K}a_{j,\infty}(x,\sigma,x) \in C^{\infty}_B(\R_+,M^{\mu_j;(\ell_1,\ell_2)}_O(Z)),
$$
and consequently
\begin{gather*}
\opm(a_{j,0}(x,\sigma,x)) : x^{\alpha}H^{s;(\ell_1,\ell_2)}_b(\R_+\times Z) \to x^{\alpha}H^{s-\mu_j;(\ell_1,\ell_2)}_b(\R_+\times Z), \\
\opm(a_{j,\infty}(x,\sigma,x)) : x^{\alpha}H^{s;(\ell_1,\ell_2)}_b(\R_+\times Z) \to x^{\alpha+K}H^{s-\mu_j;(\ell_1,\ell_2)}_b(\R_+\times Z)
\end{gather*}
for all $s,\alpha \in \R$. By the mapping properties of $G_1$ we obtain
$$
G_1 \opm(a_{2,0}(x,\sigma,x)) : x^{\alpha}H^{s;(\ell_1,\ell_2)}_b(\R_+\times Z) \to x^{\alpha'}H^{s';(\ell_1,\ell_2)}_b(\R_+\times Z)
$$
for all $s,s' \in \R$ and all $\alpha' \leq \alpha$, and
$$
G_1 \opm(a_{2,\infty}(x,\sigma,x)) : x^{\alpha}H^{s;(\ell_1,\ell_2)}_b(\R_+\times Z) \to x^{\alpha'}H^{s';(\ell_1,\ell_2)}_b(\R_+\times Z)
$$
for all $s,s' \in \R$ and all $\alpha' \leq \alpha + K$, and so also for all $\alpha' \leq \alpha$, which shows that
$$
G_1 \opm(a_{2}(x,\sigma,x)) : x^{\alpha}H^{s;(\ell_1,\ell_2)}_b(\R_+\times Z) \to x^{\alpha'}H^{s';(\ell_1,\ell_2)}_b(\R_+\times Z).
$$
for all $s,s' \in \R$ and all $\alpha' \leq \alpha$. Because
$$
G_2 : x^{\alpha}H^{s;(\ell_1,\ell_2)}_b(\R_+\times Z) \to x^{\alpha'}H^{s'+\mu_1;(\ell_1,\ell_2)}_b(\R_+\times Z)
$$
we get
$$
\opm(a_{1,0}(x,\sigma,x))G_2 : x^{\alpha}H^{s;(\ell_1,\ell_2)}_b(\R_+\times Z) \to x^{\alpha'}H^{s';(\ell_1,\ell_2)}_b(\R_+\times Z)
$$
for all $s,s'  \in \R$ and $\alpha' \leq \alpha$, and similarly
$$
G_2 : x^{\alpha}H^{s;(\ell_1,\ell_2)}_b(\R_+\times Z) \to x^{\alpha'-K}H^{s'+\mu_1;(\ell_1,\ell_2)}_b(\R_+\times Z)
$$
and
$$
\opm(a_{1,\infty}(x,\sigma,x))G_2 : x^{\alpha}H^{s;(\ell_1,\ell_2)}_b(\R_+\times Z) \to x^{\alpha'}H^{s';(\ell_1,\ell_2)}_b(\R_+\times Z),
$$
and so
$$
\opm(a_1(x,\sigma,x))G_2 : x^{\alpha}H^{s;(\ell_1,\ell_2)}_b(\R_+\times Z) \to x^{\alpha'}H^{s';(\ell_1,\ell_2)}_b(\R_+\times Z)
$$
for all $s,s' \in \R$ and all $\alpha' \leq \alpha$. For the formal adjoints we have
\begin{align*}
[G_1\opm(a_2(x,\sigma,x))]^* &= [\opm(a_2(x,\sigma,x))]^*G_1^*, \\
[\opm(a_1(x,\sigma,x))G_2]^* &= G_2^*[\opm(a_1(x,\sigma,x))]^*,
\end{align*}
and because $[\opm(a_j)(x,\sigma,x)]^* \in \Psi^{\mu_j;\vec{\ell}}_{\textup{(cl)}}(\R_+\times Z)$ by Theorem~\ref{AdjointsCalculus} we obtain with the above that both $\opm(a_1(x,\sigma,x))G_2$ and $G_1\opm(a_2(x,\sigma,x))$ are residual.

It remains to consider the composition $\opm(a_1(x,\sigma,x))\opm(a_2(x,\sigma,x))$. We first consider
$$
a_j(x,\sigma,\tau)  \in C^{\infty}_B(\R_+,M^{-\infty}_{O}(Z;\overline{\R}_+)).
$$
We can then compose $\opm(a_1(x,\sigma,x))\opm(a_2(x,\sigma,x)) = \opm(c(x,\sigma))$ by the standard Mellin pseudodifferential calculus, and the formulas \eqref{MelLeibniz} and \eqref{MelLeibnizRem} apply. For every $N \in \N$
\begin{align*}
c(x,\sigma) &= \sum\limits_{k=0}^{N-1}\frac{1}{k!}[\partial_{\sigma}^ka_1](x,\sigma,x)[(xD_x + \tau D_{\tau})^ka_2](x,\sigma,x) + r_N(x,\sigma) \\
&= c_N(x,\sigma) + r_{N}(x,\sigma).
\end{align*}
We further decompose
$$
r_N(x,\sigma) = r_{N,0}(x,\sigma) + r_{N,\infty}(x,\sigma)
$$
with
\begin{align*}
r_{N,0}(x,\sigma) = \frac{1}{2\pi}\! &\int_0^1\!\frac{(1-\theta)^{N-1}}{(N-1)!} \iint y^{-i\eta} \cdot \\
&[\partial_{\sigma}^Na_1](x,\sigma+\theta\eta,x) [(xD_{x} + \tau D_{\tau})^Na_{2,0}](xy,\sigma,xy)\,\frac{dy}{y}d\eta d\theta, \\
r_{N,\infty}(x,\sigma) = \frac{1}{2\pi}\! &\int_0^1\!\frac{(1-\theta)^{N-1}}{(N-1)!} \iint y^{-i\eta} \cdot \\
&[\partial_{\sigma}^Na_1](x,\sigma+\theta\eta,x) [(xD_{x} + \tau D_{\tau})^Na_{2,\infty}](xy,\sigma,xy)\,\frac{dy}{y}d\eta d\theta,
\end{align*}
where $a_{2,0} = \omega a_2$ and $a_{2,\infty} = (1-\omega)a_2$ as above.

\medskip

\noindent
Let $\mu_1',\mu_2' \in \R$, $j_0 \in \N_0$, and $\mu_0 < 0$ be arbitrary.

We first analyze $r_{N,0}(x,\sigma)$: For $0 \leq j \leq j_0$ write
\begin{align*}
x^jr_{N,0}(x,\sigma) = \frac{1}{2\pi}\! &\int_0^1\!\frac{(1-\theta)^{N-1}}{(N-1)!} \iint y^{-i\eta} \cdot \\
&[\tau^j\partial_{\sigma}^Na_1](x,\sigma+\theta\eta,x) [(xD_{x} + \tau D_{\tau})^Na_{2,0}](xy,\sigma,xy)\,\frac{dy}{y}d\eta d\theta.
\end{align*}
Because
$$
C^{\infty}_B(\R_+,M^{\mu_2';\vec{\ell}}_{O}(Z;\overline{\R}_+)) \ni a_2(x,\sigma,\tau) \mapsto a_{2,0}(x,\sigma,x) \in C^{\infty}_B(\R_+,M^{\mu_2';(\ell_1,\ell_2)}_{O}(Z))
$$
is continuous we see from the formula that there exists $N_0 \in \N$ such that for all $N \geq N_0$ and all $0 \leq j \leq j_0$ the maps
$$
(a_1,a_2) \mapsto x^jr_{N,0}(x,\sigma) \textup{ and } (a_1,a_2) \mapsto x^jr_{N,0}(x,\sigma-ij)
$$
are continuous in
$$
C^{\infty}_B(\R_+,M^{\mu_1';\vec{\ell}}_{O}(Z;\overline{\R}_+)) \times C^{\infty}_B(\R_+,M^{\mu_2';\vec{\ell}}_{O}(Z;\overline{\R}_+)) \to C^{\infty}_B(\R_+,M^{\mu_0;(\ell_1,\ell_2)}_{O}(Z)),
$$
and so
\begin{gather*}
\opm(r_{N,0}),\, [\opm(r_{N,0})]^* = x^{-j} [\opm(x^jr_{N,0}(x,\sigma-ij))]^* : \\
x^{\alpha}H^{s;(\ell_1,\ell_2)}_b(\R_+\times Z) \to x^{\alpha-j}H^{s-\mu_0;(\ell_1,\ell_2)}_b(\R_+\times Z)
\end{gather*}
for all $\alpha,s \in \R$.

We next want to show the same mapping properties for $\opm(r_{N,\infty})$ for $N$ sufficiently large. Pick $K > 0$ large enough such that
$$
C^{\infty}_B(\R_+,M^{\mu_2';\vec{\ell}}_{O}(Z;\overline{\R}_+)) \!\ni\! a_2(x,\sigma,\tau) \mapsto x^{-K}a_{2,\infty}(x,\sigma,x) \!\in\! C^{\infty}_B(\R_+,M^{\mu_2';(\ell_1,\ell_2)}_{O}(Z))
$$
is continuous. We use the analyticity of the operator-valued symbols in the oscillatory integral formula for $r_{N,\infty}(x,\sigma)$ to shift the integration contour and write for every $0 \leq j \leq j_0$
\begin{gather*}
x^jr_{N,\infty}(x,\sigma) = \frac{1}{2\pi}\! \int_0^1\!\frac{(1-\theta)^{N-1}}{(N-1)!} \iint y^{-i\eta} \cdot \\
[\tau^{j+K}\partial_{\sigma}^Na_1](x,\sigma+\theta(\eta-iK),x)[x^{-K}(xD_{x}+\tau D_{\tau})^Na_{2,\infty}](xy,\sigma,xy)\,\frac{dy}{y}d\eta d\theta.
\end{gather*}
This formula shows that we can find $N_0 \in \N$ such that for $N \geq N_0$ and $0 \leq j \leq j_0$ the maps
$$
(a_1,a_2) \mapsto x^jr_{N,\infty}(x,\sigma) \textup{ and } (a_1,a_2) \mapsto x^jr_{N,\infty}(x,\sigma-ij)
$$
are continuous in
$$
C^{\infty}_B(\R_+,M^{\mu_1';\vec{\ell}}_{O}(Z;\overline{\R}_+)) \times C^{\infty}_B(\R_+,M^{\mu_2';\vec{\ell}}_{O}(Z;\overline{\R}_+)) \to C^{\infty}_B(\R_+,M^{\mu_0;(\ell_1,\ell_2)}_{O}(Z)),
$$
so $\opm(r_{N,\infty})$ and $[\opm(r_{N,\infty})]^*$ have the desired mapping properties. We conclude that
$$
\opm(r_{N}),\, [\opm(r_{N})]^* : x^{\alpha}H^{s;(\ell_1,\ell_2)}_b(\R_+\times Z) \to x^{\alpha-j}H^{s-\mu_0;(\ell_1,\ell_2)}_b(\R_+\times Z)
$$
for all $\alpha,s \in \R$.

Now let $a_j(x,\sigma,\tau) \in C^{\infty}_B(\R_+,M^{\mu_j;\vec{\ell}}_{O,\textup{(cl)}}(Z;\overline{\R}_+))$. By Lemma~\ref{AnalyticSmoothingApprox} there exist sequences $a_{j,\nu} \in C^{\infty}_B(\R_+,M^{-\infty}_{O}(Z;\overline{\R}_+))$ such that $a_{j,\nu}(x,\sigma,\tau) \to a_j(x,\sigma,\tau)$ as $\nu \to \infty$ in $C^{\infty}_B(\R_+,M^{\mu_j;\vec{\ell}}_{O,\textup{(cl)}}(Z;\overline{\R}_+))$ for $\mu_j' > \mu_j$.

For every $u \in \S_0(\R_+\times Z)$ and $N \in \N$ large enough
\begin{gather*}
\opm(a_1(x,\sigma,x))\opm(a_2(x,\sigma,x))u \longleftarrow \opm(a_{1,\nu}(x,\sigma,x))\opm(a_{2,\nu}(x,\sigma,x))u \\
= \opm(c_{N,\nu})u + \opm(r_{N,\nu})u \longrightarrow \opm(c_N)u + \opm(r_N)u.
\end{gather*}
So $\opm(a_1(x,\sigma,x))\opm(a_2(x,\sigma,x)) = \opm(c_N) + \opm(r_N)$ for $N$ large enough, and $\opm(r_N)$ has the mapping properties previously shown. Consequently, if the operator-valued symbol $a_1{\#}a_2 \in C^{\infty}_B(\R_+,M^{\mu_1+\mu_2;\vec{\ell}}_{O,\textup{(cl)}}(Z;\overline{\R}_+))$ has the asymptotic expansion \eqref{LeibnizAsymptotics}, then
$$
G_0 = \opm(a_1(x,\sigma,x))\opm(a_2(x,\sigma,x)) - \opm((a_1{\#}a_2)(x,\sigma,x))
$$
satisfies
$$
G_0,\,G_0^* : x^{\alpha}H^{s;(\ell_1,\ell_2)}_b(\R_+\times Z) \to x^{\alpha'}H^{s';(\ell_1,\ell_2)}_b(\R_+\times Z)
$$
for all $\alpha,s,s' \in \R$ and all $\alpha' \leq \alpha$ by Lemma~\ref{ImprovedMappingProps} and the mapping properties of $\opm(r_N)$ (as $N \to \infty$). The theorem is proved.
\end{proof}

\end{document}
